%!TEX root = ../main.tex
\chapter{Background}
\label{cap:background}

Questo capitolo introduce, da zero, ogni concetto usato nel resto della
tesi, così da poter leggere i Capitoli~\ref{cap:tsc-modelli} e
\ref{cap:confronto} senza assumere familiarità pregressa con nessuno dei
campi coinvolti. I quattro strumenti generali discussi (reti neurali,
apprendimento per rinforzo, reti neurali ricorrenti, grafi e reti neurali a
grafo) sono sviluppati in letteratura per problemi che non hanno nulla a
che vedere con il traffico: la loro applicazione a questo problema
specifico è discussa solo nell'ultima sezione del capitolo e formalizzata
nel dettaglio nel Capitolo~\ref{cap:tsc-modelli}.

\section{Struttura del capitolo}
\label{sec:background-overview}

Il capitolo comprende cinque sezioni. Questa sezione ne anticipa il
contenuto, motivando anche perché alcuni argomenti correlati non vengano
discussi, così da poter leggere il resto del capitolo senza doverne
indovinare i confini.

Nella §\ref{sec:background-nn} si introducono i concetti minimi di rete
neurale (percettrone, strati, funzione di loss, discesa del gradiente,
backpropagation) su cui ogni rete usata in questo lavoro, dal DQN alla
LSTM alle reti a grafo, è costruita. Non si discute l'apprendimento
supervisionato in profondità, né alcuna architettura per dati a griglia
regolare come le \emph{convolutional neural network}: lo stato osservato
da un'intersezione in questo lavoro è un vettore di conteggi e tempi di
attesa per corsia, non un'immagine, per cui l'elaborazione di immagini non
è mai rilevante.

Nella §\ref{sec:background-rl} si introduce il reinforcement learning,
dal paradigma agente-ambiente fino alle estensioni del DQN effettivamente
usate in questo lavoro, passando per la distinzione tra apprendimento
model-based e model-free e per l'apprendimento multi-agente. Il
reinforcement learning model-based, pur introdotto per completezza
definitoria, non è discusso oltre la sua definizione: nessun modello
usato in questa tesi apprende o usa un modello esplicito della dinamica
dell'ambiente (§\ref{subsec:model-free}).

Nella §\ref{sec:background-rnn} si introduce la necessità di una memoria
per dati sequenziali e la cella LSTM che questo lavoro adotta per
integrare l'informazione temporale. Si discute brevemente anche la GRU,
un'alternativa più leggera, solo per contrasto: questo lavoro non la
utilizza.

Nella §\ref{sec:background-gnn} si introduce il concetto di grafo da zero
e il \emph{message passing} su cui ogni \emph{graph neural network} è
basata, con particolare attenzione alla GCN e alla GAT, le due
formulazioni confrontate sperimentalmente nel Capitolo~\ref{cap:confronto}
tramite le rispettive varianti meta-condizionate. Meccanismi di
comunicazione multi-agente alternativi alle GNN, come i protocolli di
comunicazione appresi esplicitamente tra agenti (non basati su un grafo di
adiacenza), esulano dallo scopo di questo lavoro, che adotta
sistematicamente la comunicazione locale via GNN discussa in questa
sezione, e non vengono quindi trattati.

Infine, nella §\ref{sec:background-tsc} le quattro sezioni precedenti
confluiscono nel problema applicativo di questo lavoro: si introduce il
vocabolario di base del controllo semaforico, le tre generazioni di
metodi di controllo (a tempi fissi, reattivo, appreso) e si posiziona
esplicitamente MetaSTGAT, il modello di riferimento di questo lavoro,
all'interno di questa tassonomia.

\section{Reti Neurali}
\label{sec:background-nn}

Una \textbf{rete neurale} è una funzione parametrica che trasforma un
input in un output attraverso una successione di trasformazioni semplici,
i cui parametri (i \textbf{pesi}) vengono appresi dai dati invece di
essere scritti a mano. L'unità di base è il \textbf{neurone artificiale}
(o \emph{percettrone}): riceve un vettore di input $x \in \mathbb{R}^n$,
ne calcola una combinazione lineare pesata $z = w^\top x + b$ (con $w$
vettore di pesi e $b$ un termine costante detto \emph{bias}), e applica a
$z$ una \textbf{funzione di attivazione} non lineare $\sigma$, producendo
l'output $a=\sigma(z)$. Attivazioni comuni sono la sigmoide $\sigma(z) =
1/(1+e^{-z})$, che comprime $z$ in $(0,1)$ ed è quindi utile ovunque serva
un valore interpretabile come probabilità o come "quanto lasciar passare"
(i \emph{gate} dell'LSTM, §\ref{sec:background-rnn}, ne sono un esempio),
la tangente iperbolica $\tanh$, che lo comprime invece in $(-1,1)$, e la
\emph{ReLU} $\max(0,z)$, che si limita ad azzerare i valori negativi. La
non linearità è essenziale: senza di essa, comporre più trasformazioni
lineari produrrebbe comunque una singola trasformazione lineare, vanificando
il vantaggio di avere più strati.

Una rete \emph{feedforward} (o \emph{multi-layer perceptron}, MLP)
dispone i neuroni in \textbf{strati} (\emph{layer}) successivi: l'output
di ogni strato è l'input dello strato seguente, dall'\emph{input layer}
(i dati grezzi) attraverso uno o più \emph{hidden layer} fino
all'\emph{output layer} finale. In forma matriciale, un singolo strato
calcola $h^{(l+1)} = \sigma\!\left(W^{(l)} h^{(l)} + b^{(l)}\right)$, dove
$h^{(l)}$ raccoglie gli output di tutti i neuroni dello strato $l$ e
$W^{(l)}$ è la matrice che impila i pesi di ciascun neurone di quello
strato. Impilare più strati permette alla rete di comporre trasformazioni
via via più astratte dell'input originale.

\paragraph{Addestramento: discesa del gradiente e backpropagation.} Una
rete neurale non nasce già utile: i suoi pesi partono da valori casuali e
vengono corretti iterativamente confrontando l'output prodotto con quello
desiderato. Una \textbf{funzione di loss} $\mathcal{L}(\hat y, y)$
quantifica la distanza tra l'output prodotto $\hat y$ e quello desiderato
$y$ (o, nel caso del reinforcement learning discusso a seguire, tra una
stima e un target calcolato da altre stime). L'addestramento consiste nel
regolare ogni peso nella direzione che riduce $\mathcal{L}$, un'operazione
nota come \textbf{discesa del gradiente}:
\[w \leftarrow w - \eta \frac{\partial \mathcal{L}}{\partial w},\]
dove $\eta$, il \textbf{tasso di apprendimento} (\emph{learning rate}),
controlla l'ampiezza di ogni aggiornamento. Calcolare la derivata parziale
$\partial \mathcal{L}/\partial w$ per ogni peso della rete, dal più vicino
all'output al più lontano, richiede applicare ripetutamente la regola
della catena del calcolo differenziale attraverso tutti gli strati
intermedi: l'algoritmo che esegue questo calcolo in modo efficiente, una
sola passata dall'output verso l'input, è la \textbf{backpropagation}.
Nella pratica, il gradiente si stima su un piccolo sottoinsieme casuale di
dati (\emph{mini-batch}) alla volta invece che sull'intero insieme di
dati disponibile, una variante nota come discesa del gradiente
stocastica, molto più economica computazionalmente a fronte di una stima
del gradiente più rumorosa ma sufficiente a guidare comunque
l'addestramento verso una soluzione utile.

\paragraph{Oltre la discesa del gradiente pura.} La regola di
aggiornamento vista sopra usa lo stesso tasso di apprendimento $\eta$ per
ogni peso della rete, un limite quando la superficie di loss ha curvatura
molto diversa da un peso all'altro (alcune direzioni richiederebbero un
passo grande, altre uno minimo): un ottimizzatore più sofisticato adatta
l'ampiezza dell'aggiornamento parametro per parametro, sulla base della
storia recente dei gradienti osservati. Il \textbf{momentum} accumula una
media mobile esponenziale dei gradienti passati e aggiorna i pesi in quella
direzione media invece che nella sola direzione del gradiente corrente,
smorzando le oscillazioni lungo le direzioni in cui il segno del gradiente
cambia di continuo. \textbf{RMSprop} \citep{tieleman2012rmsprop} adotta un
principio complementare: mantiene per ciascun peso una media mobile del
quadrato dei gradienti recenti e divide l'aggiornamento per la sua radice,
riducendo il passo sui pesi il cui gradiente è stato storicamente grande in
valore assoluto e aumentandolo su quelli il cui gradiente è stato piccolo,
così da rendere il progresso più uniforme tra pesi con scale molto diverse.
\textbf{Adam} \citep{kingma2015adam} combina entrambe le idee (media mobile
del gradiente e del suo quadrato) ed è oggi l'ottimizzatore di riferimento
per la maggior parte delle reti neurali profonde. \textbf{Questo lavoro
adotta RMSprop}, coerentemente con il modello di riferimento
\citep{wang2022metastgat}, per l'addestramento di ogni rete $Q$ descritta
nel Capitolo~\ref{cap:tsc-modelli}.

Ogni rete usata in questo lavoro, dal DQN (§\ref{sec:background-rl}) alla
LSTM (§\ref{sec:background-rnn}) alle reti a grafo
(§\ref{sec:background-gnn}), è in ultima analisi una particolare
composizione di questi stessi due ingredienti: strati di neuroni con
un'attivazione non lineare, addestrati per discesa del gradiente tramite
backpropagation. Ciò che cambia da un'architettura all'altra è come gli
strati sono collegati tra loro, non il principio di addestramento
sottostante.

\section{Reinforcement Learning}
\label{sec:background-rl}

\subsection{Il paradigma agente-ambiente}

Il \emph{reinforcement learning} (RL, apprendimento per rinforzo) è un
paradigma di apprendimento automatico in cui un \textbf{agente} impara a
comportarsi interagendo ripetutamente con un \textbf{ambiente}, senza che
nessuno gli mostri esempi di comportamento corretto. A ogni passo di
tempo $t$, l'agente osserva lo stato corrente dell'ambiente, sceglie
un'\textbf{azione} tra quelle disponibili, e riceve in cambio due cose:
un nuovo stato (l'ambiente è cambiato, anche per effetto dell'azione
scelta) e un numero scalare, la \textbf{ricompensa}, che indica quanto
quell'azione sia stata desiderabile in quel momento. Non esiste alcun
supervisore che dica all'agente qual è l'azione giusta da compiere:
l'unico segnale disponibile è la ricompensa, ottenuta \emph{dopo} aver
agito, e spesso in ritardo rispetto alla decisione che l'ha causata (una
fase semaforica scelta ora può congestionare un incrocio vicino solo
dopo diversi passi). Il problema che il RL affronta è per l'appunto
questo: imparare, per tentativi ripetuti, quale comportamento massimizzi
la somma delle ricompense ricevute nel tempo, non solo quella immediata.

Questa formulazione è deliberatamente generale: lo stesso paradigma
descrive un programma che impara a giocare a scacchi, un braccio robotico
che impara ad afferrare un oggetto, e, come in questa tesi, un
semaforo che impara a scegliere quale fase attivare. Ciò che cambia da un
problema all'altro sono solo le definizioni concrete di stato, azione e
ricompensa.

Vale la pena distinguere questo paradigma dall'\textbf{apprendimento
supervisionato}, il paradigma di machine learning più diffuso e spesso il
primo che si incontra: lì una rete neurale (§\ref{sec:background-nn})
riceve coppie input-output già corrette (ad esempio migliaia di immagini,
ciascuna già etichettata con il nome dell'oggetto che raffigura) e impara
a riprodurre quella corrispondenza. Nel RL non esiste alcuna etichetta
corretta da imitare: l'agente deve scoprire da solo, per tentativi, quale
azione sia buona in ciascuno stato, guidato solo da un segnale scalare di
ricompensa che oltretutto può arrivare con ritardo rispetto alla
decisione che l'ha causata. È un problema strutturalmente più difficile,
ma anche l'unico adatto a situazioni, come il controllo semaforico, in
cui nessuno conosce a priori quale sia la "decisione corretta" da
etichettare: lo si può solo osservare indirettamente dalle sue
conseguenze sul traffico.

\subsection{Processo decisionale di Markov}

La formalizzazione matematica standard di un problema di RL è il
\textbf{processo decisionale di Markov} (MDP) \citep{suttonbarto2018},
definito dalla tupla
\[\mathcal{M} = (\mathcal{S}, \mathcal{A}, P, R, \gamma)\]
dove $\mathcal{S}$ è l'insieme di tutti gli stati possibili dell'ambiente,
$\mathcal{A}$ l'insieme di tutte le azioni possibili, $P(s' \mid s,a)$ la
probabilità di transitare nello stato $s'$ avendo eseguito l'azione $a$
nello stato $s$, $R(s,a)$ la ricompensa attesa per quella stessa coppia, e
$\gamma \in [0,1)$ il \textbf{fattore di sconto}, che pondera l'importanza
delle ricompense future rispetto a quelle immediate: un $\gamma$ vicino a
$0$ rende l'agente "miope", interessato quasi solo alla ricompensa
immediata; un $\gamma$ vicino a $1$ lo rende attento anche a conseguenze
lontane nel tempo. Il nome "Markov" indica un'assunzione precisa: la
probabilità di transizione $P(s'\mid s,a)$ dipende solo dallo stato
corrente $s$, mai dalla sequenza di stati che ha portato fin lì, cosa che
richiede che lo stato $s$ contenga già tutta l'informazione rilevante per
decidere.

Un agente segue una \textbf{politica} $\pi(a \mid s)$, una distribuzione
di probabilità sulle azioni condizionata allo stato corrente; il suo
obiettivo è massimizzare il \textbf{ritorno} atteso
\[G_t = \sum_{k=0}^{\infty} \gamma^k r_{t+k}\]
ovvero la somma scontata delle ricompense future a partire dall'istante
$t$. Per ragionare su quanto sia buona una politica prima ancora di
eseguirla fino in fondo, si definiscono due funzioni valore: la
funzione valore-stato $V^\pi(s) = \mathbb{E}_\pi[G_t \mid s_t = s]$ e la
funzione valore-azione $Q^\pi(s,a) = \mathbb{E}_\pi[G_t \mid s_t=s,
a_t=a]$, che misurano rispettivamente quanto ritorno atteso si ottiene
seguendo $\pi$ da uno stato, o da una coppia stato-azione. Queste
soddisfano l'\textbf{equazione di Bellman}, che esprime il valore di uno
stato in funzione del valore atteso degli stati successivi, invece che
dover simulare l'intero futuro:
\begin{equation}
  V^\pi(s) = \sum_a \pi(a\mid s) \sum_{s'} P(s'\mid s,a) \big[R(s,a) + \gamma V^\pi(s')\big]
\end{equation}
Per una politica ottima $\pi^*$ che massimizza $V^\pi(s)$ per ogni stato
simultaneamente, vale l'analoga equazione di Bellman di ottimalità:
\begin{equation}
  Q^*(s,a) = \mathbb{E}_{s' \sim P(\cdot \mid s,a)}\left[
    R(s,a) + \gamma \max_{a'} Q^*(s', a')
  \right]
  \label{eq:bellman-optimal}
\end{equation}
Se si conoscesse $Q^*$, la politica ottima sarebbe immediata da ricavare:
scegliere sempre $\operatorname*{argmax}_a Q^*(s,a)$. Il problema del RL
si riduce quindi, in buona parte, a stimare $Q^*$ senza conoscerlo a
priori.

\subsection{Apprendimento model-based e model-free}
\label{subsec:model-free}

Esiste una distinzione fondamentale tra due famiglie di algoritmi RL, a
seconda che l'agente costruisca o meno una rappresentazione esplicita
della dinamica dell'ambiente. Un algoritmo \textbf{model-based} apprende
(o riceve) un modello approssimato di $P(s'\mid s,a)$ e $R(s,a)$, e lo usa
per pianificare, cioè per simulare mentalmente le conseguenze di più
sequenze di azioni prima di sceglierne una, senza dover necessariamente
eseguirle davvero nell'ambiente. Un algoritmo \textbf{model-free}, al
contrario, non costruisce alcun modello della dinamica: stima
direttamente $Q^*$ (o la politica) dall'esperienza osservata, senza mai
provare a prevedere esplicitamente lo stato successivo. Il vantaggio del
model-free è la semplicità e la minore dipendenza dalla qualità di un
modello che potrebbe essere difficile da apprendere accuratamente; lo
svantaggio è un uso tipicamente meno efficiente dell'esperienza raccolta,
poiché non la si riusa per simulare scenari mai osservati direttamente.

\textbf{Questo lavoro adotta un approccio interamente model-free}: né il
modello di riferimento \citep{wang2022metastgat} né le varianti
implementate in questa tesi (Capitolo~\ref{cap:tsc-modelli}) apprendono o
usano un modello esplicito della dinamica del traffico. L'agente stima
direttamente una funzione $Q$ dall'esperienza simulata in CityFlow, con la
famiglia di tecniche descritta nel resto di questa sezione.

\subsection{Q-learning e Deep Q-Network}

Il \textbf{Q-learning} tabellare approssima $Q^*$ applicando
iterativamente l'Equazione~\ref{eq:bellman-optimal} come regola di
aggiornamento, aggiustando il valore stimato $Q(s,a)$ verso il target $r +
\gamma \max_{a'} Q(s',a')$ osservato a ogni transizione. Il nome
"tabellare" indica che $Q$ è letteralmente una tabella con una riga per
ogni stato e una colonna per ogni azione: un approccio che richiede una
tabella con tutte le coppie $(s,a)$, il che risulta inapplicabile per
spazi di stato $\mathcal{S}$ continui o ad alta dimensionalità, come nel
caso del controllo semaforico, dove lo stato include conteggi di veicoli
e tempi di attesa che assumono innumerevoli combinazioni possibili
\citep{li2018drlsurvey}.

\paragraph{Un esempio numerico.} Si supponga un singolo incrocio con due
sole azioni possibili, "verde a Nord-Sud" ($a_1$) e "verde a Est-Ovest"
($a_2$), $\gamma=0.9$, e una tabella $Q$ ancora azzerata ovunque. In un
episodio di addestramento, dallo stato $s$ (coda presente su entrambe le
direzioni) l'agente esegue $a_1$ per esplorazione casuale
($\varepsilon$-greedy, discusso a breve), osserva una ricompensa $r=-2$
(qualche veicolo è comunque rimasto in coda a Est-Ovest) e transita in un
nuovo stato $s'$ in cui $\max_{a'} Q(s',a')$ vale, poniamo, $5$ (una stima
già presente da esperienza precedente). L'aggiornamento tabellare del
Q-learning corregge la stima verso il target di Bellman con un tasso
$\eta$:
\[Q(s,a_1) \leftarrow Q(s,a_1) + \eta\Big[\underbrace{r + \gamma
\max_{a'}Q(s',a')}_{\text{target} \,=\, -2 + 0.9 \cdot 5 \,=\, 2.5} -
Q(s,a_1)\Big] = 0 + \eta(2.5 - 0) = 2.5\,\eta.\]
Ripetendo questo aggiornamento su molte transizioni campionate, la stima
$Q(s,a_1)$ converge verso il vero ritorno atteso dell'eseguire $a_1$ in
$s$ e poi seguire la politica ottima in seguito, senza che nessuno abbia
mai detto esplicitamente all'agente se $a_1$ fosse la scelta giusta in
$s$: l'informazione emerge indirettamente dalla combinazione di
ricompensa osservata e stima del valore dello stato successivo. Il DQN,
descritto di seguito, sostituisce la tabella con una rete neurale ma
applica esattamente lo stesso principio di aggiornamento verso il target
di Bellman.

\paragraph{Deep Q-Network (DQN).} \citet{mnih2015dqn} superano tale limite
sostituendo la tabella con una rete neurale $Q(s,a;\theta)$, che
generalizza a stati mai visti esattamente interpolando tra stati simili
osservati durante l'addestramento, i cui parametri (pesi) $\theta$ vengono
aggiornati per discesa del gradiente in modo che $Q(s,a;\theta)$ si
avvicini al target di Bellman. Per mitigare l'instabilità
dell'addestramento dovuta alla correlazione temporale dei campioni (stati
consecutivi di uno stesso episodio sono molto simili tra loro, violando
l'assunzione di campioni indipendenti richiesta dalla discesa del
gradiente stocastico), gli autori introducono inoltre il \emph{replay
buffer}: una memoria in cui si accumulano le transizioni $(s,a,r,s')$
osservate, da cui si campionano poi mini-batch in ordine casuale per
l'addestramento, rompendo la correlazione temporale.

\subsection{Esplorazione $\varepsilon$-greedy}

Un agente che seguisse sempre e solo l'azione a valore $Q$ stimato più
alto non raccoglierebbe mai esperienza sulle conseguenze di azioni
alternative, rischiando di convergere su una politica sub-ottima per via
di una mancata esplorazione: se per puro caso, nelle prime fasi
dell'addestramento, un'azione appare leggermente migliore di un'altra, un
agente completamente "goloso" (\emph{greedy}) non la riprova mai più, e
non scopre mai se l'altra fosse in realtà preferibile. La strategia più
comune per bilanciare sfruttamento (\emph{exploitation}, scegliere
l'azione ritenuta migliore) ed esplorazione (\emph{exploration},
raccogliere informazione su azioni meno note) è $\varepsilon$-greedy: a
ogni passo di decisione, con probabilità $\varepsilon$ l'agente sceglie
un'azione uniformemente a caso, e con probabilità $1-\varepsilon$ sceglie
$\operatorname*{argmax}_a Q(s,a;\theta)$. Il valore di $\varepsilon$
tipicamente decresce nel corso dell'addestramento, da un valore iniziale
vicino a 1 (quasi solo esplorazione, quando la stima di $Q$ è ancora
inaffidabile) verso un valore finale piccolo ma non nullo (quasi solo
sfruttamento, mantenendo comunque un minimo di esplorazione residua).

\subsection{Estensioni usate in questo lavoro}

Il DQN di base ammette diverse estensioni indipendenti, ciascuna delle
quali corregge una specifica fonte di instabilità o inefficienza; la
procedura di training descritta nel Capitolo~\ref{cap:tsc-modelli} le
adotta tutte, motivandole caso per caso rispetto all'articolo di
riferimento:
\begin{itemize}
  \item Il \textbf{Double DQN} \citep{vanhasselt2016double} disaccoppia la
    selezione dell'azione dalla sua valutazione, usando la rete online per
    scegliere l'azione migliore e una \emph{target network}
    $Q(\cdot;\theta^-)$, una copia periodicamente sincronizzata della rete
    online, solo per valutarla:
    $y = r + \gamma\, Q\!\left(s', \operatorname*{argmax}_{a'}
    Q(s',a';\theta);\ \theta^-\right)$. Questo meccanismo evita il problema
    del \emph{moving target}, che si verifica ottimizzando una stima
    rispetto a parametri in continua variazione, e riduce il bias
    sistematico di sovrastima del DQN standard (usare la stessa rete sia
    per scegliere sia per valutare l'azione tende a preferire azioni il
    cui valore è per caso sovrastimato dal rumore di stima).
  \item Il \textbf{Prioritized Experience Replay} (PER,
    \citealp{schaul2016per}) sostituisce il campionamento uniforme dal
    replay buffer con un campionamento proporzionale all'errore di
    differenza temporale (\emph{TD error}) $\delta_i$, la discrepanza tra
    il valore stimato e il target di Bellman per quella transizione: le
    transizioni che la rete predice peggio vengono ricampionate più
    spesso, concentrando la capacità di apprendimento dove serve di più.
    Assegnando a ciascuna transizione una priorità $p_i = |\delta_i| +
    \varepsilon$, la probabilità di campionamento è $P(i) = p_i^\alpha /
    \sum_k p_k^\alpha$, con $\alpha \in [0,1]$ iperparametro che ne
    controlla la dipendenza dall'errore. Per correggere il bias introdotto
    da un campionamento non uniforme, la loss viene pesata tramite
    coefficienti di \emph{importance sampling} $w_i = \left(N \cdot
    P(i)\right)^{-\beta}$, dove $N$ è la dimensione del buffer e $\beta \in
    [0,1]$ iperparametro dinamico che ne controlla il grado di correzione.
  \item La \textbf{Huber loss} (o \emph{smooth L1}) sostituisce l'errore
    quadratico medio (MSE) con una funzione quadratica vicino allo zero e
    lineare lontano da esso, risultando meno sensibile agli outlier
    prodotti da transizioni con errore di stima molto grande (con MSE, un
    singolo errore enorme domina il gradiente dell'intero batch; con Huber
    il suo contributo cresce solo linearmente):
    \begin{equation}
      L_\delta(y, \hat y) =
      \begin{cases}
        \tfrac{1}{2}(y-\hat y)^2 & \text{se } |y-\hat y| \le \delta \\
        \delta\left(|y-\hat y| - \tfrac{1}{2}\delta\right) & \text{altrimenti}
      \end{cases}
    \end{equation}
    Tipicamente $\delta = 1$.
  \item Il \textbf{soft update} (\emph{Polyak averaging}) sostituisce
    l'aggiornamento periodico discreto della \emph{target network} con una
    media mobile esponenziale: $\theta^- \leftarrow \tau\,\theta +
    (1-\tau)\,\theta^-$ con $\tau \ll 1$, stabilizzando le variazioni del
    target a ogni passo di ottimizzazione invece di farlo saltare
    bruscamente ogni volta che la copia periodica scatta.
  \item Quando la rete $Q$ include un modulo ricorrente
    (§\ref{sec:background-rnn}), l'assunzione di Markov sulle singole
    transizioni campionate dal replay buffer viene invalidata: lo stato
    ricorrente $(h,c)$ usato durante il training viene inizializzato a
    zero, introducendo un disallineamento rispetto allo stato generato
    sequenzialmente durante la raccolta dati, un problema noto come
    \emph{hidden state staleness}. Seguendo \citet{kapturowski2019r2d2}, si
    campionano intere sequenze di lunghezza $L$ e si dedicano i primi passi
    della sequenza (\emph{burn-in}) esclusivamente a ricostruire uno stato
    ricorrente plausibile, senza calcolare gradiente, prima di eseguire la
    \emph{backpropagation through time} (BPTT, §\ref{sec:background-rnn})
    sui passi restanti.
\end{itemize}

\subsection{Apprendimento multi-agente}
\label{subsec:marl}

Tutto quanto descritto finora assume un solo agente che interagisce con
un solo ambiente. Quando più agenti operano nello stesso ambiente,
ciascuno osservando solo una porzione dello stato globale e influenzando
con le proprie azioni anche ciò che gli altri agenti osservano e
ricevono come ricompensa, il problema si generalizza da MDP a
\textbf{processo decisionale di Markov decentralizzato e parzialmente
osservabile} (Dec-POMDP): "decentralizzato" perché non esiste un'unica
azione congiunta scelta da un solo decisore, ma $N$ azioni scelte
indipendentemente da $N$ agenti; "parzialmente osservabile" perché ogni
agente vede solo il proprio stato locale, non lo stato globale
completo dell'ambiente. Quest'ultima proprietà non è solo teorica: in
questo lavoro si concretizza nel campo visivo limitato di ogni
intersezione (Capitolo~\ref{cap:tsc-modelli}), che esclude deliberatamente
dallo stato osservato ciò che accade oltre una certa distanza dalla
corsia, ben prima che la questione multi-agente entri in gioco.

Il multi-agente introduce una difficoltà che il caso a singolo agente non
ha: la \textbf{non stazionarietà}. L'equazione di Bellman e l'intera teoria
del Q-learning assumono che l'ambiente risponda in modo coerente nel
tempo alla stessa azione nello stesso stato; ma se più agenti apprendono
simultaneamente, ciascuno sta anche cambiando il proprio comportamento
mentre gli altri lo osservano, cosicché per ogni singolo agente
l'ambiente (che include il comportamento degli altri agenti) appare
cambiare regole nel tempo, anche se le vere dinamiche fisiche sottostanti
restano identiche. Nessuna delle garanzie di convergenza valide per il RL
a singolo agente si applica quindi automaticamente al caso multi-agente.

Una seconda difficoltà, distinta dalla non stazionarietà, è
l'\textbf{assegnazione del merito} (\emph{credit assignment}): quando più
agenti contribuiscono a un solo risultato osservato, non è ovvio quale
frazione di quel risultato attribuire a ciascuno. Il caso classico è una
singola ricompensa globale condivisa da tutti gli agenti: un agente non ha
modo diretto di distinguere quanto di quella ricompensa dipenda dalla
propria azione e quanto dalle azioni altrui, un problema che soluzioni
come COMA \citep{li2018drlsurvey} affrontano stimando un termine
controfattuale, quanto sarebbe cambiata la ricompensa se il singolo agente
avesse agito diversamente a parità di comportamento altrui. Il problema di
questo lavoro è più mite: ogni intersezione riceve una ricompensa $r_i^t$
calcolata esclusivamente sul proprio stato locale (Capitolo~%
\ref{cap:tsc-modelli}), non una ricompensa globale condivisa, evitando così
la versione più severa del problema appena descritto. Una versione
attenuata resta comunque presente: lo stato locale di un'intersezione, e
quindi la sua stessa ricompensa, dipende in parte dalle code di traffico
generate dalle scelte delle intersezioni a monte, per cui
l'effetto della propria azione e quello dell'azione altrui restano in
parte intrecciati anche con una ricompensa puramente locale.

Due schemi di coordinamento ricorrono in letteratura
\citep{li2018drlsurvey}: l'\emph{addestramento centralizzato con
esecuzione decentralizzata} (CTDE), in cui gli agenti condividono
informazione globale solo durante il training per poi decidere
localmente, e l'\emph{addestramento ed esecuzione entrambi
decentralizzati} (DTDE), in cui l'informazione resta sempre e solo
locale, anche durante l'addestramento. Una scelta indipendente da questa
è se i diversi agenti condividano o meno gli stessi pesi (\emph{parameter
sharing}): condividerli riduce il numero di parametri totali e permette a
un agente di beneficiare dell'esperienza raccolta da tutti gli altri, a
costo di una minore capacità di specializzarsi su condizioni locali
peculiari, a meno che il meccanismo non trovi un altro modo per
differenziare il comportamento tra agenti pur condividendo i pesi.

\textbf{Questo lavoro è un sistema multi-agente}: ogni intersezione della
rete stradale è un agente a sé, con un proprio stato $s_i^t$, una propria
azione $a_i^t$ (la fase semaforica scelta) e una propria ricompensa
$r_i^t$, calcolati indipendentemente per ogni intersezione $i$
(Capitolo~\ref{cap:tsc-modelli}); è per questo che il titolo di questa
tesi parla esplicitamente di apprendimento per rinforzo
\emph{multi-agente}. L'architettura è di tipo DTDE, con scambio di
informazione locale tra intersezioni vicine (§\ref{sec:background-gnn}), e
adotta il parameter sharing: tutte le intersezioni condividono la stessa
rete $Q$, ma un meccanismo di meta-learning (§\ref{sec:background-gnn})
genera pesi diversi per nodi con caratteristiche locali diverse, così da
attenuare proprio il costo della condivisione appena descritto.

\section{Reti Neurali Ricorrenti}
\label{sec:background-rnn}

Una rete neurale \emph{feedforward} calcola il proprio output come
funzione del solo input corrente: ogni esempio è trattato in modo
indipendente dai precedenti, e la rete non conserva alcuna memoria di
cosa abbia visto in passato. Questa assunzione fallisce ogni volta che i
dati hanno una struttura sequenziale in cui l'ordine e la storia contano
quanto il singolo valore osservato (testo, audio, serie temporali, o lo
stato di un'intersezione stradale nei passi decisionali precedenti)
\citep{lipton2015rnnreview}. Le \emph{reti neurali ricorrenti} (RNN)
affrontano il problema introducendo un ciclo: a ogni istante $t$, un nodo
riceve in ingresso non solo il dato corrente $x^{(t)}$ ma anche il proprio
stato nascosto al passo precedente $h^{(t-1)}$,
\[h^{(t)} = \sigma\!\left(W^{hx} x^{(t)} + W^{hh} h^{(t-1)} + b_h\right),\]
così che l'informazione rilevante di tutta la sequenza osservata fino a
$t$ possa, in linea di principio, persistere in $h^{(t)}$: lo stesso
insieme di pesi $W^{hx}, W^{hh}, b_h$ viene riapplicato a ogni istante di
tempo, invece di avere pesi diversi per ogni passo. "Srotolando" la rete
lungo l'asse del tempo, cioè disegnando esplicitamente una copia dei nodi
per ogni istante $t=1,2,\dots$ collegata alla successiva, ogni passo
diventa uno strato di una rete feedforward molto profonda a pesi
condivisi, addestrabile con la stessa backpropagation vista in precedenza
applicata lungo la sequenza anziché lungo i layer: questo algoritmo
prende il nome di \emph{backpropagation through time} (BPTT).

\paragraph{Il problema del gradiente.} Una rete ricorrente addestrata su
sequenze lunghe soffre di un problema strutturale: il gradiente
retropropagato attraverso molti passi temporali tende a annullarsi
(\emph{vanishing gradient}) o a divergere (\emph{exploding gradient}) a
seconda che i pesi ricorrenti abbiano modulo minore o maggiore di uno,
poiché il gradiente si ottiene moltiplicando ripetutamente per lo stesso
fattore a ogni passo srotolato, esattamente come una potenza $c^n$ tende a
zero o esplode al crescere di $n$ a seconda che $|c|$ sia minore o
maggiore di uno. Nella pratica, questo rende di fatto impossibile
apprendere dipendenze che coprono più di poche decine di passi con una
RNN semplice.

\paragraph{Long Short-Term Memory (LSTM).} L'LSTM \citep{lipton2015rnnreview}
risolve il problema sostituendo il semplice nodo ricorrente con una
\emph{cella di memoria}: uno stato interno $c^{(t)}$ collegato a se stesso
da un percorso a peso fisso unitario (il \emph{constant error carousel}),
attraverso cui il gradiente può fluire per molti passi senza annullarsi né
esplodere, non dovendo più moltiplicare ripetutamente per un peso appreso
potenzialmente minore o maggiore di uno. L'accesso a questo stato è
regolato da tre \emph{gate}, unità sigmoidali il cui valore in $[0,1]$
modula quanta informazione passi attraverso di esse, in modo analogo a una
valvola che si apre o si chiude in proporzione al segnale che riceve. Con
pesi $W_g \in \mathbb{R}^{2d_1 \times d_1}$ e bias $b_g \in
\mathbb{R}^{d_1}$ condivisi da ogni nodo, per ciascuno dei quattro gate
$g \in \{\text{forget, input, output, cell}\}$:
\begin{equation}
  \begin{bmatrix} f_i^t \\ \iota_i^t \\ o_i^t \\ g_i^t \end{bmatrix} =
  \begin{bmatrix} \sigma \\ \sigma \\ \sigma \\ \tanh \end{bmatrix}
  \left(
  W_{\{f,\iota,o,g\}}^{\top} [\,e_i^t \,\|\, h_i^{t-1}\,] + b_{\{f,\iota,o,g\}}
  \right),
  \label{eq:lstm-background-gates}
\end{equation}
\[ c_i^t = f_i^t \odot c_i^{t-1} + \iota_i^t \odot g_i^t, \]
\[ x_i^t = h_i^t = o_i^t \odot \tanh(c_i^t); \]
il \emph{forget gate} $f$ decide quanto dello stato precedente
dimenticare, l'\emph{input gate} $\iota$ quanto del nuovo contenuto
candidato $g$ (talvolta indicato $\tilde c$ in altre notazioni) incorporare
nello stato aggiornato, e l'\emph{output gate} $o$ quanto dello stato
risultante esporre come output della cella. Un valore di forget gate
vicino a 1 lascia passare l'informazione pregressa quasi inalterata anche
dopo molti passi; uno vicino a 0 la cancella, permettendo alla rete di
"dimenticare" deliberatamente informazione non più rilevante. La
notazione con pedice $i$ e apice $t$, anziché quella generica appena
introdotta, anticipa l'uso di questa stessa equazione
(Equazione~\ref{eq:stgat-lstm-gates}) nel Capitolo~\ref{cap:tsc-modelli}
per il modulo ricorrente di MetaSTGAT, dove $i$ indicizza l'intersezione
ed $e_i^t$ sostituisce l'input generico $x^{(t)}$: ogni intersezione
mantiene quindi una propria cella di memoria, aggiornata a ogni passo
decisionale con il proprio stato osservato.

\paragraph{Gated Recurrent Unit (GRU).} Una variante più leggera dell'LSTM,
la \textbf{GRU} \citep{cho2014gru}, fonde il \emph{forget gate} e
l'\emph{input gate} in un unico \emph{update gate} e rinuncia allo stato di
cella $c^{(t)}$ separato, esponendo un solo stato nascosto $h^{(t)}$ sia
come memoria interna sia come output: meno parametri da addestrare e un
grafo computazionale più semplice, a fronte di una capacità di
rappresentazione leggermente inferiore rispetto ai quattro gate distinti
dell'LSTM. Le due varianti raggiungono nella pratica prestazioni simili su
molti compiti, e la scelta tra loro è spesso guidata più dalla
convenzione di un'architettura di riferimento che da un vantaggio netto
dell'una sull'altra: questo lavoro adotta l'LSTM per coerenza con il
modello di riferimento \citep{wang2022metastgat}.

\section{Grafi e Graph Neural Network}
\label{sec:background-gnn}

\subsection{Cos'è un grafo}

Un \textbf{grafo} $G = (V, E)$ è una struttura matematica composta da un
insieme di \textbf{nodi} (o vertici) $V$ e un insieme di \textbf{archi}
$E$, dove ogni arco collega una coppia di nodi. È lo strumento standard
per rappresentare entità e le relazioni tra loro, quando quelle relazioni
non hanno una struttura regolare a griglia (come i pixel di un'immagine) né
sequenziale (come le parole di una frase): una rete di amicizie tra
persone, le pagine web collegate da link, le molecole i cui atomi sono
legati da legami chimici, e, in questo lavoro, una rete stradale, dove i
nodi sono le intersezioni e gli archi le strade che le collegano
direttamente.

Un arco può essere \textbf{diretto} (indica un verso, come "$u$ segue $v$"
sui social network) o \textbf{non diretto} (indica solo una relazione
reciproca, come "$u$ è amico di $v$"); può inoltre avere un \textbf{peso}
associato, un numero che quantifica l'intensità o il costo della
relazione (ad esempio la lunghezza di una strada). L'insieme dei nodi
collegati direttamente a un nodo $i$ da un arco si chiama
\textbf{vicinato} di $i$, indicato $\mathcal{N}(i)$; il numero di archi
incidenti su $i$ è il suo \textbf{grado}. Una sequenza di nodi collegati
uno al successivo da un arco è un \textbf{cammino}; la lunghezza del
cammino più breve tra due nodi qualunque è la loro \textbf{distanza} sul
grafo, misurata in numero di archi, non in senso fisico o geografico:
due intersezioni collegate da una singola strada molto lunga sono a
distanza $1$ tanto quanto due collegate da una strada brevissima. Questa
nozione di distanza torna utile più avanti, quando si discuterà quanti
passi di propagazione servano perché l'informazione di un nodo raggiunga
un altro nodo della rete. Un grafo con $I$ nodi si può
rappresentare in forma matriciale con la \textbf{matrice di adiacenza} $A
\in \{0,1\}^{I \times I}$, dove $A_{iu}=1$ se e solo se esiste un arco tra
$i$ e $u$ (o il peso dell'arco, se il grafo è pesato, invece del semplice
$1$); $A$ è simmetrica se e solo se il grafo non è diretto. È spesso utile
aggiungere a ogni nodo un \textbf{self-loop}, un arco da un nodo verso se
stesso ($A_{ii}=1$), cosicché un nodo sia incluso nel proprio stesso
vicinato: la notazione $\hat A = A + I$ (con $I$ qui la matrice identità,
da non confondere con il numero di nodi) indica questa aggiunta e ricorre
nelle equazioni seguenti.

Un piccolo esempio concreto: tre intersezioni disposte in fila, $1-2-3$,
dove $2$ è collegata sia a $1$ sia a $3$, ma $1$ e $3$ non sono collegate
direttamente. Il vicinato di $2$ è $\mathcal{N}(2) = \{1,3\}$, il suo
grado è $2$; i vicinati di $1$ e $3$ sono invece $\mathcal{N}(1)=\{2\}$ e
$\mathcal{N}(3)=\{2\}$, entrambi di grado $1$. La matrice di adiacenza
corrispondente, non diretta e senza pesi, è
\[A = \begin{pmatrix} 0&1&0 \\ 1&0&1 \\ 0&1&0 \end{pmatrix},\]
simmetrica perché il grafo non è diretto: se $2$ è vicino di $1$, anche
$1$ è vicino di $2$. Con l'aggiunta dei self-loop, $\hat A = A+I$ ha
inoltre un $1$ su ogni elemento della diagonale.

Nella rete stradale di questo lavoro, ogni nodo $i$ del grafo porta con sé
un vettore di caratteristiche (o \emph{embedding}) che ne descrive lo
stato corrente, ad esempio quanti veicoli attendono su ciascuna corsia:
il problema che le GNN risolvono è come far sì che l'embedding di un nodo
tenga conto anche di ciò che accade nei nodi a lui collegati, non solo del
proprio stato isolato.

\subsection{Message passing}

Le \emph{Graph Neural Network} (GNN) affrontano questo problema con un
principio comune, il \textbf{message passing} \citep{wu2021gnnsurvey,
bacciu2020gentlegnn}. A ogni iterazione (o \emph{layer}) $l$, ogni nodo
$i$ con rappresentazione corrente $h_i^{(l)}$ esegue due operazioni in
successione: prima un passo di \textbf{aggregazione}, che combina i
messaggi ricevuti da ciascun vicino $u \in \mathcal{N}(i)$ in un unico
vettore di messaggio $m_i$,
\[m_i = \mathrm{AGGREGATE}\big(\{h_u^{(l)} : u \in \mathcal{N}(i)\}\big),\]
poi un passo di \textbf{aggiornamento}, che combina il messaggio
aggregato con la rappresentazione precedente dello stesso nodo per
produrre quella nuova,
\[h_i^{(l+1)} = \mathrm{UPDATE}\big(h_i^{(l)},\ m_i\big).\]
Tutte le GNN discusse in questa tesi, incluse le due formulazioni di
seguito, sono istanze di questo stesso schema generale, che differiscono
solo nella scelta concreta di AGGREGATE e UPDATE: nella GCN, AGGREGATE è
una media pesata dalla sola topologia; nella GAT, un peso appreso in base
al contenuto dei due nodi coinvolti. Impilando più layer, l'informazione
può viaggiare oltre il solo vicinato immediato: dopo $k$ layer, la
rappresentazione di un nodo dipende (indirettamente) da tutti i nodi
raggiungibili entro $k$ archi di distanza, poiché il messaggio ricevuto da
un vicino al layer $l$ incorpora già, ricorsivamente, i messaggi che
quel vicino aveva ricevuto dai propri stessi vicini al layer $l-1$.

\paragraph{Graph Convolutional Network (GCN).} \citet{kipf2017gcn}
propongono la formulazione più basilare di message passing tra quelle qui
discusse: la rappresentazione di ciascun nodo si aggiorna mediando quelle
dei vicini con un peso fisso, determinato unicamente dalla topologia
locale (quanti vicini ha un nodo), non dal contenuto delle
rappresentazioni stesse. Indicando con $\hat A = A + I$ la matrice di
adiacenza con self-loop e con $\hat D$ la matrice diagonale dei gradi
corrispondente ($\hat D_{ii} = \sum_u \hat A_{iu}$), la regola di
propagazione di un singolo layer si scrive, in forma matriciale:
\begin{equation}
  H^{(l+1)} = \sigma\!\left(\hat D^{-1/2} \hat A\, \hat D^{-1/2} H^{(l)} W^{(l)}\right)
  \label{eq:gcn-propagation}
\end{equation}
dove $H^{(l)} \in \mathbb{R}^{I \times d_l}$ raccoglie le rappresentazioni
di tutti gli $I$ nodi del grafo al layer $l$, $W^{(l)}$ è la matrice di
pesi condivisa da ogni nodo e $\sigma$ una non linearità puntuale. Per un
singolo nodo $i$, la normalizzazione simmetrica $\hat D^{-1/2}\hat
A\hat D^{-1/2}$ equivale a pesare il contributo di ogni vicino $u$ con
$1/\sqrt{\hat d_i \hat d_u}$:
\begin{equation}
  h_i' = \sigma\!\left(\sum_{u \in \mathcal{N}(i) \cup \{i\}}
    \frac{1}{\sqrt{\hat d_i \hat d_u}}\, W h_u\right)
  \label{eq:gcn-aggregate}
\end{equation}
un peso identico per ogni arco del vicinato di $i$, funzione solo del
grado dei due nodi coinvolti e mai del loro contenuto. Tornando all'esempio
a tre nodi $1-2-3$ di §\ref{sec:background-gnn}, i gradi con self-loop sono
$\hat d_1 = \hat d_3 = 2$ e $\hat d_2 = 3$: il nodo centrale $2$ aggrega il
proprio vicinato $\{1,2,3\}$ pesando ciascun contributo rispettivamente
$1/\sqrt{\hat d_2 \hat d_1} = 1/\sqrt{6}$, $1/\sqrt{\hat d_2 \hat d_2} =
1/3$ e $1/\sqrt{\hat d_2 \hat d_3} = 1/\sqrt{6}$: due pesi identici verso i
vicini $1$ e $3$, indipendentemente da quale dei due porti, ad esempio, il
segnale di una coda di veicoli più lunga. È, per costruzione,
il termine di paragone più naturale rispetto a cui misurare cosa aggiunge
un meccanismo di attenzione appreso come quello discusso di seguito: il
Capitolo~\ref{cap:tsc-modelli} confronta sperimentalmente proprio questa
formulazione (MetaSTGCN) con l'attenzione appresa del modello di
riferimento.

\paragraph{Graph Attention Network (GAT).} \citet{velickovic2018gat}
sostituiscono il peso fisso della GCN con un coefficiente di attenzione
appreso, specifico per ciascuna coppia di nodi collegati: il contributo di
un vicino non dipende più soltanto dalla topologia, ma anche dal suo
contenuto. Dato un nodo $i$ con vicinato $\mathcal{N}(i)$, embedding di
input $h_i$, matrice di proiezione condivisa $W$ e vettore di attenzione
$\vec a$:
\begin{equation}
  e_{iu} = \mathrm{LeakyReLU}\!\left(\vec a^\top [W h_i \,\|\, W h_u]\right),
  \qquad
  \alpha_{iu} = \frac{\exp(e_{iu})}{\sum_{u' \in \mathcal{N}(i)} \exp(e_{iu'})}
  \label{eq:gat-attention}
\end{equation}
\begin{equation}
  h_i' = \sigma\!\left(\sum_{u \in \mathcal{N}(i)} \alpha_{iu}\, W h_u\right)
  \label{eq:gat-aggregate}
\end{equation}
dove $\|$ denota la concatenazione: $e_{iu}$ è un punteggio di
compatibilità grezzo tra $i$ e $u$, che il \emph{softmax} sul vicinato
trasforma in un coefficiente $\alpha_{iu} \in [0,1]$ che somma a uno su
tutti i vicini di $i$, esattamente come una distribuzione di probabilità.
Con $K$ teste di attenzione indipendenti (\emph{multi-head attention}),
ciascuna con i propri pesi $W$ e $\vec a$, gli output si concatenano nei
layer intermedi e si mediano nell'ultimo, secondo lo schema introdotto per
l'attenzione da \citet{vaswani2017attention}: teste diverse possono così
imparare a specializzarsi su aspetti diversi della relazione tra due nodi.
A differenza della GCN, il coefficiente $\alpha_{iu}$ consente a un nodo
di ponderare diversamente i propri vicini in base al loro contenuto, non
solo alla loro posizione nel grafo: una proprietà rilevante per
un'intersezione i cui vicini abbiano importanza marcatamente diversa,
come nel caso di un'arteria principale contrapposta a una strada
secondaria. La funzione di compatibilità additiva appena descritta, per
quanto la più nota, non è l'unica possibile: il Capitolo~\ref{cap:tsc-modelli}
(Meta-GAT) ne impiega una variante strutturalmente diversa, con Query, Key
e Value esplicite e distinte al posto della proiezione $W$ condivisa, e un
prodotto scalare al posto della concatenazione.

\subsection{Invarianza e equivarianza per permutazione}
\label{subsec:permutation-invariance}

Sia l'aggregazione della GCN sia quella della GAT condividono una
proprietà che vale la pena rendere esplicita, perché è ciò che permette a
queste reti di funzionare su grafi di dimensione diversa senza dover
riaddestrare nulla. Entrambe calcolano il messaggio aggregato $m_i$ di un
nodo $i$ come una somma (pesata, uniformemente nella GCN o per attenzione
nella GAT) sui vicini $u \in \mathcal{N}(i)$: una somma non dipende
dall'ordine in cui i suoi addendi vengono elencati, quindi $m_i$ non
dipende dall'ordine in cui i vicini di $i$ vengono numerati o elencati, ma
solo da quali essi siano. Questa proprietà si chiama \textbf{invarianza per
permutazione} dell'operatore di aggregazione, ed è ciò che distingue
strutturalmente una GNN da una rete feedforward applicata a un vettore di
caratteristiche di lunghezza fissa: una MLP a cui si desse in input, in un
qualche ordine arbitrario, le caratteristiche di ciascun nodo del grafo
concatenate assegnerebbe implicitamente un significato specifico a ciascuna
posizione del vettore (il primo blocco di ingressi, il secondo, ...) e
produrrebbe quindi un output diverso se semplicemente si rinumerassero i
nodi, pur descrivendo esattamente la stessa rete stradale;
\citet{wei2019colight} osservano esattamente questo limite per gli schemi
di coordinamento basati su una lista ordinata di vicini
(§\ref{sec:background-tsc}). Applicando la stessa proprietà a ogni nodo del
grafo simultaneamente, l'intero layer di message passing è
\textbf{equivariante per permutazione}: rinumerare i nodi del grafo in
ingresso produce in uscita esattamente le stesse rappresentazioni,
solo rinumerate allo stesso modo, invece di un risultato arbitrariamente
diverso \citep{bronstein2021geometric}.

Questa proprietà ha una conseguenza pratica diretta per la generalizzazione
discussa nel Capitolo~\ref{cap:confronto}: le matrici di peso $W$ di un
layer GCN o GAT non fanno alcuna assunzione sul numero di nodi del grafo
$I$ o sul grado di un nodo specifico, poiché la stessa operazione di
aggregazione, con gli stessi pesi, si applica indipendentemente al vicinato
di ciascun nodo, quale che sia la sua cardinalità. Una rete addestrata
su una griglia $4\times4$ (16 nodi) può quindi essere applicata, senza
alcuna modifica ai pesi né riaddestramento, a una griglia $5\times5$ o
$6\times6$ (25 o 36 nodi): l'unico requisito è che il vicinato di un
incrocio in una griglia più grande assomigli, dal punto di vista della
distribuzione delle caratteristiche osservate, a quello visto in
addestramento. È esattamente questa proprietà, non un meccanismo aggiunto
appositamente, a rendere possibile in linea di principio l'intera
valutazione di generalizzazione a topologie di rete più grandi che il
Capitolo~\ref{cap:confronto} conduce su MetaSTGAT e sulle sue varianti.

\subsection{Reti spatio-temporali su grafo}

Una rete stradale non è soltanto un grafo: è un grafo il cui stato
(occupazione delle corsie, fase attiva) evolve nel tempo, e le due
dimensioni risultano intrecciate, poiché il valore osservato a un nodo in
un dato istante è correlato sia con i valori passati allo stesso nodo sia
con quelli ai nodi limitrofi. La famiglia di modelli che combina un layer
di propagazione sul grafo (GCN o GAT, appena visti) con una componente
ricorrente (LSTM, §\ref{sec:background-rnn}) sull'asse temporale prende il
nome di \emph{spatiotemporal graph neural network} (STGNN), un termine
ormai consolidato nella letteratura sul forecasting di serie temporali
multivariate correlate \citep{cini2025graphdl}. MetaSTGAT, descritto nel
dettaglio nel Capitolo~\ref{cap:tsc-modelli}, costituisce un esempio di
questa famiglia applicata al controllo semaforico anziché alla previsione:
un layer di \emph{graph attention} cattura la dipendenza spaziale tra
intersezioni limitrofe, un layer ricorrente (LSTM) ne cattura l'evoluzione
temporale, e i due vengono addestrati congiuntamente end-to-end.

\subsection{Pesi generati da meta-learning}

Sia la GCN sia la GAT, nella loro formulazione originale, condividono gli
stessi pesi su tutti i nodi del grafo (esattamente il \emph{parameter
sharing} introdotto in §\ref{subsec:marl}): un'assunzione ragionevole
quando i nodi sono sufficientemente omogenei, meno quando, come in una
rete stradale reale, la topologia locale di ciascun nodo (numero di
strade entranti, presenza di un'arteria, posizione periferica o centrale)
può differire sensibilmente da nodo a nodo. Un \emph{hypernetwork}
affronta il problema generando i pesi di un layer non come parametri
fissi appresi direttamente, bensì come output di una seconda rete,
condizionata sulle caratteristiche locali del nodo: è di fatto una via di
mezzo tra pesi completamente condivisi e pesi completamente indipendenti
per ogni nodo, che riduce il costo del parameter sharing senza rinunciare
del tutto ai suoi benefici. Il Capitolo~\ref{cap:tsc-modelli} descrive nel
dettaglio come MetaSTGAT applichi questo principio sia al layer di
\emph{graph attention} sia al layer ricorrente.

\subsection{Propagazione a lungo raggio}

Sia la GCN sia la GAT propagano informazione a un solo passo di distanza
(un arco nel grafo) per layer: raggiungere un vicino a distanza $k$
richiede impilarne altrettanti, con un costo di parametri crescente e un
rischio noto di \emph{over-smoothing}, per il quale le rappresentazioni
dei nodi divengono troppo omogenee al crescere della profondità (ogni nodo
finisce per "vedere" praticamente tutto il grafo, perdendo la capacità di
distinguersi dai propri vicini). Meccanismi più recenti, tra cui
\textbf{SONAR} (\citealp{trenta2025sonar}), affrontano il problema della
propagazione a lungo raggio secondo un principio diverso, ispirato a
sistemi dinamici oscillatori piuttosto che alla concatenazione sequenziale
di layer: il dettaglio di questo meccanismo e la sua applicazione al
controllo semaforico sono rimandati al Capitolo~\ref{cap:tsc-modelli}.

Un esempio numerico chiarisce perché il numero di passi $L$ di
propagazione (il parametro che, in un meccanismo come SONAR, gioca il
ruolo che altrove è giocato dal numero di layer impilati) sia una scelta
tutt'altro che neutra. In una griglia stradale come quella $4\times4$
usata per l'addestramento in questo lavoro, un'intersezione interna ha 4
vicini diretti (Nord, Sud, Est, Ovest): a distanza $1$ nel senso definito
in §\ref{sec:background-gnn} si trovano quindi fino a 4 nodi. A distanza
$2$ si aggiungono i vicini dei vicini: fino a 8 nodi ulteriori (i 4 vicini
diagonali più altre 4 intersezioni allineate a due passi), per un totale
di 12. In generale, su una griglia regolare il numero di nodi raggiungibili
entro $L$ passi cresce quadraticamente con $L$, non linearmente: passare da
$L=2$ a $L=4$ non raddoppia il campo recettivo di un'intersezione, lo
espande di un fattore prossimo a 4. Un campo recettivo più ampio non è però
un vantaggio gratuito. Ogni passo aggiuntivo di propagazione richiede alla
rete di retropropagare il gradiente attraverso uno stadio computazionale in
più, e aumenta il rischio di \emph{over-smoothing} discusso sopra: le
rappresentazioni dei nodi coinvolti in un raggio via via più ampio tendono
a somigliarsi, rendendo più difficile per il modello distinguere
un'intersezione congestionata da una vicina che non lo è. Il Capitolo~%
\ref{cap:confronto} riprende esattamente questo compromesso, confrontando
empiricamente MetaSTSONAR con $L=2$ e $L=4$.

\section{Traffic Signal Control come problema}
\label{sec:background-tsc}

\subsection{Vocabolario di base}

Il controllo semaforico (\emph{Traffic Signal Control}, TSC) governa il
transito dei veicoli in un'intersezione stradale. Il vocabolario di base,
qui ripreso da \citet{wei2020tscsurvey}, è il seguente. Un'
\textbf{approach} (approccio) è una via che porta all'intersezione,
composta da una o più corsie in ingresso e, tipicamente, altrettante in
uscita. Un \textbf{movimento} è il passaggio di un veicolo da una corsia
in ingresso a una specifica corsia in uscita (ad esempio: prosecuzione
dritto, svolta a destra o a sinistra), e a ogni movimento è associato un
\textbf{segnale}, verde se il movimento è autorizzato, rosso altrimenti.
Non tutti i movimenti possono ricevere il verde simultaneamente: due
movimenti le cui traiettorie si intersecano fisicamente sono in
\textbf{conflitto}, e la loro compatibilità reciproca si registra in una
\emph{matrice dei conflitti}. Un insieme di movimenti a due a due
compatibili, che possono quindi ricevere il verde nello stesso istante,
definisce una \textbf{fase}; il Capitolo~\ref{cap:tsc-modelli} riporta le
otto fasi usate in questo lavoro su un'intersezione a quattro bracci. Una
\textbf{sequenza di fasi} stabilisce l'ordine in cui le fasi si succedono
nel tempo, e un \textbf{piano semaforico} ne fissa anche la durata; quando
la sequenza si ripete identica nel tempo si parla di piano
\textbf{ciclico}, con \textbf{ciclo} la durata di una ripetizione completa.

\subsection{Tre generazioni di controllo}

Il TSC si articola in tre approcci, in ordine crescente di complessità e,
storicamente, di introduzione \citep{shams2018atsctaxonomy}.

\textbf{Controllo a tempi fissi.} Un piano semaforico ciclico, con fasi di
durata predeterminata calcolata su stime di traffico medie, si ripete
senza mai osservare lo stato reale della rete. È un approccio robusto, in
quanto non può fallire per un sensore guasto, ma per costruzione non
reattivo: la stessa sequenza viene eseguita indipendentemente da
congestioni anomale, incidenti o variazioni del traffico non modellate in
fase di pianificazione. Metodi classici dell'ingegneria dei trasporti,
come quello di Webster per una singola intersezione isolata, calcolano
analiticamente la durata di ciclo e la ripartizione delle fasi che
minimizza il tempo di attesa atteso sotto un'assunzione di flusso
stazionario \citep{wei2020tscsurvey}, un'ipotesi che nella pratica vale
raramente per l'intera giornata. Il metodo di Webster, ad esempio, calcola
la durata di ciclo desiderata come
\[C_0 = \frac{1.5L + 5}{1 - Y},\]
dove $L$ è il tempo perso complessivo per ciclo (la somma dei brevi
intervalli in cui nessun movimento riceve il verde, ad esempio durante il
cambio di fase) e $Y$ è la somma, su tutte le fasi, del rapporto tra flusso
osservato e flusso di saturazione della fase più critica: un ciclo più
lungo riduce l'incidenza percentuale del tempo perso ma allunga l'attesa
media di chi deve aspettare il proprio turno, il compromesso centrale di
ogni piano a tempi fissi. Il piano risultante resta valido solo finché il
traffico reale resta vicino ai volumi $Y$ usati per calcolarlo.

\textbf{Controllo reattivo.} Una regola locale osserva lo stato corrente
dell'intersezione (tipicamente il numero di veicoli in coda per corsia) e
sceglie l'azione che ottimizza un criterio istantaneo, senza affidarsi a
un piano prestabilito. All'interno di questa categoria,
\citet{shams2018atsctaxonomy} distinguono ulteriormente il grado di
adattività: il controllo \emph{attuato} più semplice si limita a
prolungare o abbreviare le fasi di un piano di base ancora sostanzialmente
predefinito (estendere una fase finché un sensore rileva veicoli sulla
corsia servita, ad esempio), mentre un controllo pienamente
\emph{adattivo} ricalcola la struttura stessa del piano, non solo la sua
durata, in risposta al traffico osservato. Piattaforme come SCATS e SCOOT
si collocano in questa seconda fascia, coordinando intere aree urbane
aggiustando periodicamente i piani semaforici sulla base del traffico
osservato, pur restando ancorate a una struttura di piano pre-esistente
\citep{wei2020tscsurvey}. \textsc{Max-Pressure}
\citep{varaiya2013maxpressure} rappresenta invece una rottura netta con
questa impostazione: modellando il movimento dei veicoli come una rete di
code \emph{store-and-forward} (ogni corsia è una coda che riceve veicoli
da monte e li inoltra a valle), definisce la \textbf{pressione} di un
movimento $(r_i \to r_o)$, dalla corsia in ingresso $r_i$ alla corsia in
uscita $r_o$, come la differenza tra il numero di veicoli in coda sulle
due corsie, $p(r_i \to r_o) = n(r_i) - n(r_o)$, dove $n(\cdot)$ conta i
veicoli presenti sulla corsia: un valore positivo indica che la corsia a
monte è più congestionata di quella a valle, quindi che servire quel
movimento libererebbe più coda di quanta se ne creerebbe. La pressione di
un'intera fase $\phi$, che autorizza più movimenti compatibili
contemporaneamente (§\ref{sec:background-tsc}), è la somma delle pressioni
dei movimenti che quella fase abilita, $P(\phi) = \sum_{(r_i \to r_o) \in
\phi} p(r_i \to r_o)$; l'algoritmo sceglie a ogni istante la fase
$\operatorname*{argmax}_\phi P(\phi)$, senza alcun piano né ciclo
predefinito. Il risultato teorico
centrale dell'articolo originale è che Max-Pressure è \emph{massimamente
stabilizzante}: se è possibile stabilizzare la rete, cioè non far
divergere la lunghezza delle code nel tempo, Max-Pressure la stabilizza,
usando solo la lunghezza delle code sulle corsie adiacenti
all'intersezione e senza conoscere a priori il volume di domanda in
ingresso. Max-Pressure è tuttavia miope per costruzione: non pianifica
alcun compromesso su un orizzonte temporale non immediato, e la condizione
di stabilità non tiene in considerazione le attese tra le diverse
direzioni servite da un'intersezione, un aspetto di equità che questo
lavoro esamina esplicitamente (Capitolo~\ref{cap:confronto}).

\textbf{Controllo appreso.} L'approccio basato sul reinforcement learning
(§\ref{sec:background-rl}) consente l'ottimizzazione diretta di metriche
di rete globali (es. \emph{travel time} o percentuale di completamento)
sostituendo euristiche come la pressione istantanea, al netto del costo
computazionale dell'addestramento simulato. Applicato a una rete di più
intersezioni, il problema si formalizza come un Dec-POMDP
(§\ref{subsec:marl}): ogni intersezione osserva solo il proprio stato
locale, ma le sue azioni influenzano, con un ritardo proporzionale alla
distanza spaziale, lo stato delle intersezioni adiacenti alterando il
flusso di veicoli che le attraversano. \textbf{Questo lavoro si colloca
interamente in questa terza categoria}: ogni intersezione è un agente RL
model-free (§\ref{subsec:model-free}) che decide quale fase attivare
osservando solo il proprio stato locale.

\paragraph{Architetture di coordinamento.} Esistono tre strategie distinte
che affrontano la dipendenza tra intersezioni vicine. Un agente
centralizzato che decide le azioni di tutta la rete con un'unica politica
elimina il problema per costruzione, ma lo spazio delle azioni congiunte
cresce esponenzialmente con il numero di intersezioni ed è impraticabile
oltre poche decine di nodi. All'estremo opposto, un agente indipendente
per intersezione, che non condivide alcuna informazione con i propri
vicini: questo approccio scala senza difficoltà ma tratta ogni
intersezione come se fosse isolata, ignorando che le sue azioni si
propagano alla rete circostante. La via intermedia, adottata anche da
questo lavoro, fa comunicare agenti indipendenti condividendo informazione
locale (architettura DTDE, §\ref{subsec:marl}). La difficoltà, discussa da
\citet{wei2019colight}, risiede nel fatto che la trasmissione dello stato
dei vicini in una sequenza predefinita (una lista ordinata di vicini)
impone a tutte le intersezioni della rete gli stessi pesi interpretativi
per ciascuna posizione della lista, senza tenere conto della reale
importanza che ogni specifico nodo adiacente può rivestire per le diverse
intersezioni del sistema. Una GNN (§\ref{sec:background-gnn}) evita il
problema: pesando il contributo di ogni vicino in base al suo contenuto
invece che alla sua posizione in un elenco, permette a un unico insieme di
pesi condiviso di generalizzare tra intersezioni con topologie locali
diverse, come quelle della rete adottata in questo lavoro
(Capitolo~\ref{cap:tsc-modelli}).

\paragraph{Coordinamento oltre l'adiacenza fisica.} Il grafo usato da
MetaSTGAT e da questo lavoro coincide con l'adiacenza fisica della rete
stradale: un'intersezione comunica solo con le intersezioni immediatamente
collegate da una strada. Dalla letteratura si osserva che interazioni tra
sole intersezioni adiacenti non siano sempre sufficienti a catturare tutte
le dinamiche di traffico, e si propongono diverse soluzioni per ovviare a
questo limite. Questo lavoro non affronta il tema direttamente,
esplorandolo solo in parte nel Capitolo~\ref{cap:confronto} attraverso un
meccanismo di propagazione a lungo raggio (SONAR,
§\ref{sec:background-gnn}) applicato comunque al solo grafo di adiacenza
fisica. Il problema resta comunque aperto, e framework che combinano
\emph{graph attention}, una componente ricorrente sul tempo e
reinforcement learning per il controllo semaforico continuano a essere
oggetto di ricerca attiva.

\paragraph{Una tassonomia delle soluzioni recenti.} La letteratura più
recente sul reinforcement learning applicato al TSC classifica i metodi
lungo assi in parte indipendenti da quelli discussi finora. Un primo asse
distingue, all'interno del multi-agente, le architetture CTDE e DTDE
appena introdotte. Un secondo asse raccoglie le tecniche di apprendimento
integrate nel ciclo RL, ciascuna introdotta per rispondere a un limite
specifico \citep{xiao2025tscreview, wei2020tscsurvey}: reti neurali a
grafo per ottenere una migliore coordinazione tra intersezioni, componenti
ricorrenti per catturare la dipendenza tra passi successivi a livello
temporale, transfer-learning per generalizzare rapidamente a scenari
nuovi, meta-learning per adattarsi a condizioni operative diverse, reti
generative avversarie per compensare dati di stato mancanti o incompleti.
Nessuna di queste tecniche risolve da sola il problema: la letteratura le
introduce quasi sempre in risposta a un limite preciso.

\paragraph{Dove si colloca MetaSTGAT.} Il modello di questo lavoro è la
combinazione di alcune delle tecniche appena citate. La sua è
un'architettura multi-agente DTDE con scambio locale di informazione
tramite GNN, parameter sharing attenuato dal meta-learning
(§\ref{sec:background-gnn}), un componente ricorrente (LSTM) per
catturare la dipendenza temporale oltre a quella spaziale
(collocandosi nella famiglia delle reti spatio-temporali su grafo,
STGNN), e un algoritmo di apprendimento model-free interamente basato
sulla famiglia DQN (§\ref{sec:background-rl}).

\subsection{Simulazione}

Valutare una politica di controllo semaforico su una rete reale non è
praticabile durante lo sviluppo: cambiare i tempi di un semaforo cittadino
migliaia di volte per tentativi, come richiederebbe l'addestramento di un
agente RL, non è ovviamente proponibile. Serve quindi un simulatore che
riproduca la dinamica del traffico con un dettaglio sufficiente, pur
mantenendo un costo computazionale accettabile per l'addestramento di
centinaia di episodi.

I simulatori di traffico si collocano su uno spettro di dettaglio a tre
livelli. Un simulatore \textbf{macroscopico} rappresenta il traffico come
un flusso continuo, descritto da grandezze aggregate come densità e
velocità media di un tratto stradale, senza modellare i singoli veicoli:
adatto a reti molto estese, ma troppo grossolano per catturare l'effetto di
una singola decisione semaforica su una coda specifica. Un simulatore
\textbf{microscopico} rappresenta invece ogni veicolo individualmente, con
una propria posizione, velocità e traiettoria, aggiornate a ogni istante da
un \textbf{modello di car-following} (che determina come un veicolo regola
la propria velocità in base a quello che lo precede, tipicamente per
evitare tamponamenti mantenendo una distanza di sicurezza) e da un modello
di cambio di corsia: è il livello di dettaglio più fedele, ma anche il più
costoso computazionalmente, poiché il numero di calcoli cresce con il
numero di veicoli simulati, non con la sola dimensione della rete. Un
simulatore \textbf{mesoscopico} si colloca a metà strada, modellando
gruppi di veicoli o code aggregate per corsia invece del singolo veicolo,
un compromesso diffuso proprio nel controllo semaforico basato su
Max-Pressure (§\ref{sec:background-tsc}), dove il modello store-and-forward
descritto sopra è esso stesso una forma di rappresentazione mesoscopica
delle code.

\textsc{CityFlow} \citep{zhang2019cityflow} è il
simulatore scelto in questo lavoro: è nominalmente microscopico (ogni
veicolo ha una propria posizione e velocità), ma, rispetto ad alternative
più diffuse come SUMO \citep{lopez2018sumo}, semplifica il modello di
car-following e di cambio di corsia per privilegiare la velocità di
simulazione rispetto al realismo del singolo veicolo, un compromesso
adatto quando l'obiettivo primario è addestrare un agente su un numero
elevato di episodi piuttosto che riprodurre con precisione millimetrica la
dinamica di un incrocio specifico. Il Capitolo~\ref{cap:tsc-modelli}
descrive nel dettaglio come questo lavoro configura l'ambiente CityFlow.

\bigskip
\noindent In sintesi, i quattro strumenti generali di questo capitolo
compongono il modello di questo lavoro come segue: una rete neurale
(§\ref{sec:background-nn}) stima la funzione $Q$ di un agente RL
model-free (§\ref{sec:background-rl}), uno per intersezione, in un
sistema multi-agente coordinato tramite lo scambio locale di informazione
proprio delle GNN (§\ref{sec:background-gnn}); una cella LSTM
(§\ref{sec:background-rnn}) integra questa informazione spaziale con la
sua evoluzione nel tempo; e il tutto è applicato al problema del
controllo semaforico appena introdotto. Il
Capitolo~\ref{cap:tsc-modelli} riprende ciascuno di questi ingredienti
uno a uno, componendoli nell'architettura completa.

Un filo conduttore attraversa più strumenti di questo capitolo, ed è utile
renderlo esplicito: quello del \textbf{bias induttivo}, l'insieme di
assunzioni che un modello incorpora prima ancora di osservare un solo
dato, e che ne guidano la generalizzazione verso soluzioni plausibili
invece che verso una qualunque delle infinite funzioni compatibili con
l'esperienza osservata. Questo lavoro introduce almeno due bias induttivi
espliciti, entrambi motivati dal problema applicativo piuttosto che
scelti in astratto: l'equivarianza per permutazione della GNN
(§\ref{subsec:permutation-invariance}), che assume che il comportamento
di un'intersezione non debba dipendere dall'ordine arbitrario in cui i
suoi vicini vengono elencati, ed è la base della generalizzazione a
griglie di dimensione diversa discussa nel Capitolo~\ref{cap:confronto};
e il termine di ricompensa dedicato all'attesa massima
(Capitolo~\ref{cap:tsc-modelli}), che assume che nessuna direzione di
traffico debba essere sacrificata indefinitamente, un'assunzione non
implicita in una ricompensa che ottimizzasse il solo tempo di percorrenza
aggregato. In entrambi i casi, il bias non è un limite imposto per
semplificare il problema, ma una scelta di design che riflette
direttamente uno degli obiettivi dichiarati di questo lavoro.
